\subsection{Human-Robot Interaction (HRI)}

\subsubsection{Speech.} The main communication interface with users is through audio. Speech-to-Text (STT) uses an adaptation of the WhisperLive~\cite{whisperlive} implementation with a faster-whisper backend, deploying \texttt{distil-large-v3}. The pipeline incorporates phrase-aware hotword handling and configurable language and transcription tasks. Speaker detection is complemented by a MediaPipe-based visual talking detection system, while DeepFilterNet3 provides noise cancellation. Text-to-Speech (TTS) uses the open-weight Kokoro-82M model~\cite{hexgrad_2025}, with playback timeout handling and support for headless operation. These services run in containers outside ROS and communicate through gRPC, preserving the microservice architecture.

\subsubsection{Display.} A secondary user interface is displayed on a screen in the robot's base. It provides insights into FRIDA's processes, allowing users to better understand how and when to interact with it, such as prompts to speak, live STT transcriptions, TTS dialog, detected keywords and vision detections. It is a native PyQt application that subscribes to ROS~2 topics directly and stores the dialog, logs and detection snapshots of each task for later review.

\subsubsection{Command Interpreter.} This pipeline translates natural language into robot-executable tasks. The process is handled by a fine-tuned Llama 3.1 8B model trained on a synthetically generated dataset. It is locally hosted via the Ollama framework on an Nvidia Orin AGX which provides sufficient VRAM.

The dataset is generated through a structured pipeline that begins with a set of predefined command templates which are populated with random variable entities (e.g., names, locations, and objects) and follow-up instructions. Entries are then generated by outputting a paraphrased natural-language instruction and a corresponding structured action plan composed of predefined primitive actions. During inference, the BAML framework parses the LLM's output into a strict JSON format avoiding action-level hallucinations, and a RAG system grounds commands ensuring entities are present in the current context.

    
This open source research was published by RoBorregos winning 2nd Best Student Paper at the 24th Mexican International Conference on Artificial Intelligence~\cite{10.1007/978-3-032-09037-9_24}.

\subsubsection{Multi-Command Scheduling.}\label{sec:scheduling} The GPSR test awards 200 bonus points for receiving several commands at once, so FRIDA merges them into a single plan instead of executing them one after another. Each action plan produced by the Command Interpreter is split into segments that begin with a navigation action, and a Held-Karp dynamic program~\cite{held1962dynamic} searches for the segment order with the shortest total travel. Two constraints keep every order executable: the actions of each command keep their original order, and, since FRIDA carries one object at a time, a segment from another command may only be interleaved between a pick and its delivery if it does not use the gripper. Travel cost is currently the straight-line distance between locations. The merged plan runs as a behavior tree with per-action timeouts and retries; if it fails, the remaining commands are executed sequentially, skipping actions already completed. In a simulated apartment, merging four commands reduced driven distance by 39\% and task time by 26\% compared to sequential execution.
